\eventsection{INTRODU\c{C}\~AO}

A mal\'aria permanece entre os principais desafios de sa\'ude p\'ublica da Amaz\^onia brasileira, onde fatores ambientais, clim\'aticos e territoriais ampliam a complexidade da vigil\^ancia epidemiol\'ogica \cite{arisco2025impact}. No estado do Amazonas, a distribui\c{c}\~ao espacial dispersa da popula\c{c}\~ao e as limita\c{c}\~oes log\'isticas dificultam o reconhecimento r\'apido de mudan\c{c}as no padr\~ao de transmiss\~ao, o que refor\c{c}a a relev\^ancia de ferramentas preditivas orientadas por dados.

Nos \'ultimos anos, a literatura tem mostrado que t\'ecnicas de aprendizado de m\'aquina podem contribuir para a previs\~ao de casos de mal\'aria quando combinadas a informa\c{c}\~oes epidemiol\'ogicas, clim\'aticas e espaciais \cite{monteiro2025integrating}. Entretanto, do ponto de vista operacional, a estimativa do n\'umero exato de casos nem sempre \'e suficiente. Em cen\'arios de vigil\^ancia, muitas vezes a necessidade central \'e identificar com anteced\^encia per\'iodos de risco elevado que exijam refor\c{c}o de monitoramento, investiga\c{c}\~ao local ou mobiliza\c{c}\~ao das equipes de sa\'ude.

Esse problema pode ser formulado como classifica\c{c}\~ao bin\'aria de eventos raros, situa\c{c}\~ao em que a distribui\c{c}\~ao entre classes \'e desbalanceada e a acur\'acia deixa de ser uma medida adequada de desempenho \cite{5128907}. Nesses contextos, curvas ROC e Precision-Recall, al\'em de m\'etricas como F1 e \textit{Brier Score}, tendem a oferecer leitura mais informativa da utilidade pr\'atica do modelo \cite{saito2015precision,davis2006relationship}.

Neste contexto, o objetivo deste artigo \'e comparar algoritmos representativos de \textit{bagging} e \textit{boosting} para a detec\c{c}\~ao de surtos de mal\'aria no Amazonas. Para isso, s\~ao avaliados os modelos Random Forest, XGBoost e LightGBM em um protocolo com separa\c{c}\~ao temporal estrita, atributos derivados de s\'eries temporais e an\'alise complementar de estabilidade dos resultados. A contribui\c{c}\~ao principal do trabalho est\'a em organizar uma compara\c{c}\~ao metodologicamente consistente entre essas fam\'ilias de modelos em um problema relevante para a sa\'ude na Amaz\^onia.
